% Part: computability
% Chapter: computability-theory
% Section: complete-ce-sets

\documentclass[../../../include/open-logic-section]{subfiles}

\begin{document}

\olfileid[pt]{cmp}{thy}{cce}
\olsection{Conjuntos computavelmente enumeráveis completos}

\begin{defn}
Um conjunto $A$ é um \emph{conjunto !!{computably enumerable} completo}
(sob redutibilidade muitos-para-um) se
\begin{enumerate}
\item $A$ é computavelmente enumerável, e
\item para qualquer outro conjunto computavelmente enumerável $B$,
  $B \leq_m A$.
\end{enumerate}
\end{defn}

Em outras palavras, os conjuntos computavelmente enumeráveis completos
são os conjuntos computavelmente enumeráveis ``mais difíceis''
possíveis. Eles permitem responder a perguntas sobre
\emph{qualquer} conjunto computavelmente enumerável.

\begin{thm}
$K$, $K_0$ e $K_1$ são todos conjuntos computavelmente enumeráveis
completos.
\end{thm}

\begin{proof}
Para ver que $K_0$ é completo, seja $B$ um conjunto computavelmente
enumerável qualquer. Então, para algum índice $e$,
\[
B = W_e = \Setabs{x}{\cfind{e}(x) \fdefined}.
\]
Seja $f$ a função $f(x) = \tuple{e, x}$. Então, para todo número
natural $x$, $x \in B$ se, e somente se, $f(x) \in K_0$. Em outras
palavras, $f$ reduz $B$ a~$K_0$.

Para ver que $K_1$ é completo, observe que, na demonstração de
\olref[k1]{prop:k1}, reduzimos $K_0$ a ele. Portanto, por
\olref[ppr]{prop:trans-red}, qualquer conjunto computavelmente
enumerável também pode ser reduzido a~$K_1$.

$K$ pode ser reduzido a $K_0$ de maneira muito semelhante.
\end{proof}

\begin{prob}
Dê uma redução de $K$ a $K_0$.
\end{prob}

\begin{digress}
Assim, verifica-se que todos os exemplos de conjuntos computavelmente
enumeráveis que consideramos até agora são computáveis ou completos.
Isso deve parecer estranho!{} Há exemplos de conjuntos computavelmente
enumeráveis que não são nem computáveis nem completos? A resposta é
sim, mas isso só foi estabelecido em meados da década de 1950,
independentemente, por Friedberg e Muchnik.
\end{digress}

\end{document}


